\begin{corollary}[Uniform allocation maximizes finite-budget coverage]
\label{cor:coupon-uniform-coverage}
Fix correctness $P\in(0,1]$ and $m\ge2$ possible correct keys.
Let $D_K(P,q)$ count distinct verified keys in $K$ independent draws with
key probabilities $Pq_1,\ldots,Pq_m$ and failure probability $1-P$.
For $u=(1/m,\ldots,1/m)$, every $K\ge1$ and $r=1,\ldots,m$ satisfy
\begin{equation}
 \Pr[D_K(P,u)\ge r]\ge\Pr[D_K(P,q)\ge r].
 \label{eq:coupon-uniform-coverage}
\end{equation}
Thus uniform conditional allocation maximizes \texttt{distinct@}$K$ at
fixed correctness, and more strongly maximizes every coverage tail
probability. This is a consequence of
\citet[Theorem~3, v1]{anceaume2015coupon}.
\end{corollary}
\begin{proof}
Map each correct key to a coupon of mass $Pq_c$ and every failure to the
null coupon of mass $1-P$. The cited theorem makes the time to collect $r$
distinct non-null coupons stochastically smallest when each has mass $P/m$.
The event that this time is at most $K$ is exactly $D_K\ge r$.
Zero coordinates and $P=1$ follow by continuity of finite-horizon
probabilities. Summing these tail inequalities gives the expectation claim.
\end{proof}
This sampling optimum does not assert that a training algorithm reaches the
uniform allocation.
